%% appendix/A-auxiliary-results.tex
%% Appendix A: Auxiliary Results
%% Hosts results we cite in the main text but do not want to prove inline.
%% Modeled on Joseph Genzer's Appendix A.

\section{Auxiliary Results}\label{app:auxiliary}

\todo{Content to fill in:

A.1 Borel-Cantelli Lemmas.
- First Borel-Cantelli: if \(\sum_n \Prob(A_n) < \infty\), then
  \(\Prob(A_n \text{ i.o.}) = 0\).
- Second Borel-Cantelli (independence version): if events are independent
  and \(\sum_n \Prob(A_n) = \infty\), then \(\Prob(A_n \text{ i.o.}) = 1\).
- Used in Section~\ref{subsec:convergence} for the convergence-in-probability
  but not-almost-surely counterexample.

A.2 Slutsky's theorem.
- Statement and a one-paragraph proof or reference to
  \citet[Thm.~5.5.17]{Casella2002}.

A.3 Continuous mapping theorem.
- Statement: if \(X_n \convd X\) and \(g\) is continuous on a set of
  \(\Prob_X\)-probability one, then \(g(X_n) \convd g(X)\).

A.4 Multivariate delta method.
- Statement and reference; used in Sections~\ref{subsec:scm-lowdim}
  and~\ref{subsec:fluctuations}.

A.5 Characteristic functions and the L\'evy continuity theorem.
\label{app:characteristic-functions}
- Definition: \(\varphi_X(t) = \E[e^{itX}]\); always exists.
- Examples: \(\Normal{0, 1}\) gives \(e^{-t^2/2}\), Cauchy gives \(e^{-|t|}\).
- L\'evy: \(X_n \convd X \iff \varphi_{X_n}(t) \to \varphi_X(t)\) pointwise,
  with continuity at 0 sufficient for the converse.
- Proof of the classical CLT via \((1 + t/n)^n \to e^t\) expansion.}

\clearpage
